% =====================================================================
% Abstract , final, 2026-08-13.
%
% Four quantitative facts only: the screening design, the top-ranked
% configuration, the strict encoding-axis confirmation, and the SOP reduction.
% Deliberately absent, because they belong in Results/Discussion and would
% crowd the abstract: the CTU-13 fold-1 FAR failure, unseen-attack detection,
% the CIC-IDS2017 ceiling, and T99.
%
% Removed from the submitted version: any claim that latency "consistently
% outperforms" the alternatives, and that LeakyParallel/latency is "the fastest
% variant at inference". Neither survives the strict evaluation.
% =====================================================================

\begin{abstract}
Network intrusion detection is a natural application for spiking neural
networks, whose event-driven computation suits resource-constrained deployment,
but their design space (the choice of neuron model and spike-encoding scheme)
remains poorly characterised for this task. We address that gap with a
controlled design-space study conducted entirely within the \texttt{snnTorch}
ecosystem, pairing nine neuron models with three spike encodings to form
\FNScrVariants{} configurations, evaluated on four benchmark datasets
(NSL-KDD, KDDCup99, CIC-IDS2017 and CTU-13) over five random seeds. The study is
deliberately intra-spiking: it asks which neuron and encoding a spiking detector
should use, not whether spiking detectors outperform conventional classifiers.

Ranking rules were fixed in a pre-registration before the analysis, and results
are reported under two protocol generations: a screening sweep over all
\FNScrVariants{} configurations, and a leakage-resistant confirmation that fits
every learned transform on training data alone, freezes each partition
independently of the model seed, and respects capture and scenario structure.
Under screening, \texttt{LeakyParallel} with latency encoding is the top-ranked
configuration, though \FNScrWithinCD{} configurations fall within one Nemenyi
critical difference of it. Under strict confirmation latency remains the
highest-ranked encoding, attaining a mean rank of \FNEncRank{} and the highest
terminal macro-F1 in \FNEncWins{} of \FNEncBlocks{} paired blocks.

Both preferences are heterogeneous rather than universal: on the most difficult
scenario-disjoint CTU-13 fold the preferred encoding inverts, and latency
reaches its terminal operating point earlier than rate coding on only two of
five strict protocols. Its most consistent advantage is computational: rate
encoding requires $\FNSopRatioMin\times$ to $\FNSopRatioMax\times$ more synaptic
operations per sample across every strict protocol. That separation has a closed
form. Between $\FNZeroPctMin\%$ and $\FNZeroPctMax\%$ of features are exactly
zero after preprocessing and neither code pays for them, so the input-spike
ratio is fixed by how the codes represent the active remainder; the resulting
prediction matches measurement to within $\FNPredErrHead\%$ on all five
protocols. Strict evaluation also
lowers previously optimistic estimates, most sharply on CTU-13, where a
scenario-disjoint partition with causal host features reduces a perfect score to
$\FNCtuStrict$. These results characterise where the recommended configuration
holds and where it does not.
\end{abstract}
